\section{Công đoạn 1: Mô hình ngôn ngữ lớn và phương pháp thích ứng y khoa}
\label{sec:chap2_llm_survey}

Mô hình ngôn ngữ lớn (Large Language Model -- LLM) là thành phần suy luận trung tâm của hệ thống hỏi--đáp: tiếp nhận bệnh án, trích xuất thực thể lâm sàng và tổng hợp câu trả lời. Mục này khảo sát hai nhóm mô hình: (i) mô hình đa dụng, gồm mô hình mã nguồn đóng và mô hình trọng số mở; và (ii) mô hình y khoa chuyên biệt, chia thành năm nhóm theo phương pháp thích ứng. Mỗi mô hình được trình bày theo cùng một trình tự: nền tảng và dữ liệu, phương pháp, kết quả, và cuối cùng là hạn chế.

%==========================================================
\subsection{Phạm vi khảo sát và benchmark đối chuẩn}

Khảo sát chỉ xét tác vụ thuần văn bản: đầu vào là ghi chép bệnh án ($\mathcal{M}_{\mathrm{clin}}$), bản tường trình chẩn đoán hình ảnh dạng văn bản ($\mathcal{M}_{\mathrm{img}}$), bảng xét nghiệm ($\mathcal{L}$) và y văn. Khi so sánh, khóa luận tuân theo ba quy ước. Thứ nhất, \emph{mô hình lõi} (ví dụ CHEESE-32B) được phân biệt với \emph{hệ thống ứng dụng} (ví dụ OrthoPilot, gồm agent, công cụ và đồ thị tri thức), và hai loại này không được so sánh ngang hàng. Thứ hai, điểm số trên các benchmark khác nhau không được quy đổi cho nhau vì mỗi benchmark có giao thức đánh giá riêng. Thứ ba, chuỗi suy luận (Chain-of-Thought -- CoT) mà mô hình hiển thị không được coi là giải thích trung thực (\emph{faithful}) cho quyết định của nó.

Bảng~\ref{tab:llm_benchmarks} liệt kê các benchmark được dùng để đối chuẩn trong mục này.

\begin{table}[H]
\begin{singlespace}
\centering
\fontsize{8pt}{10pt}\selectfont
\setlength{\tabcolsep}{3.5pt}
\renewcommand{\arraystretch}{1.1}
\begin{tabularx}{\textwidth}{|>{\raggedright\arraybackslash}p{3.2cm}|>{\raggedright\arraybackslash}X|>{\raggedright\arraybackslash}p{3.6cm}|}
\hline
\textbf{Benchmark} & \textbf{Nội dung} & \textbf{Độ đo chính} \\
\hline
VM14K \cite{vm14k_2025} & 14.000 câu hỏi y khoa tiếng Việt, 31 chuyên ngành & Pass@1 \\
\hline
Spanish MIR \cite{spanish_mir_2026} & 1.552 câu hỏi từ 8 kỳ thi nội trú Tây Ban Nha & Accuracy, ECE, AUROC \\
\hline
Mind the Gap \cite{mind_the_gap_2026} & Đối chuẩn các kỹ thuật ước lượng độ bất định trên mô hình suy luận & AUROC, ECE \\
\hline
EHR-Bench, EHRSHOT & Suy luận và dự đoán rủi ro trên bệnh án dọc thời gian (MIMIC-IV, Stanford) & F1 \\
\hline
OrthoBench \cite{orthopilot_2026} & 5.905 bệnh nhân cơ xương khớp, 1.000 mã bệnh, kiểm chứng ngoài tại 60 trung tâm & Accuracy chẩn đoán, quản lý điều trị \\
\hline
HealthBench & Hội thoại y khoa, chấm theo tiêu chí của hội đồng bác sĩ & Tỷ lệ phát biểu bị bác bỏ / bất định \\
\hline
\end{tabularx}
\caption{Các benchmark dùng để đối chuẩn mô hình ngôn ngữ lớn}
\label{tab:llm_benchmarks}
\end{singlespace}
\end{table}

%==========================================================
\subsection{Mô hình ngôn ngữ lớn đa dụng}

\subsubsection{Mô hình mã nguồn đóng}
Các mô hình thương mại như dòng GPT của OpenAI, Claude của Anthropic và Gemini của Google được dùng làm mốc tham chiếu hiệu năng và làm giám khảo (\emph{LLM-as-a-judge}). Tiêu biểu là GPT-4o \cite{achiam2023gpt4}: mô hình đạt $72.74\%$ trên VM14K \cite{vm14k_2025}, nhưng chỉ đạt F1 $0.3155$ trên EHR-Bench khi phân tích bệnh án dọc thời gian mà không có công cụ hỗ trợ. Các phiên bản mới hơn tiếp tục cải thiện, ví dụ GPT-5.1 đạt $67.2\%$ trên OrthoBench \cite{orthopilot_2026}. Tuy nhiên, nhóm mô hình này chỉ truy cập được qua API đám mây: dữ liệu bệnh nhân phải gửi ra máy chủ bên ngoài, trái với yêu cầu bảo mật dữ liệu y tế, và không thể truy cập logit hay trạng thái ẩn để đo độ bất định. Vì vậy, khóa luận chỉ dùng chúng làm mốc so sánh, không dùng làm mô hình nền.

\subsubsection{Mô hình trọng số mở}
Mô hình trọng số mở cho phép triển khai trong mạng nội bộ bệnh viện. Khóa luận khảo sát ba họ mô hình tiêu biểu.

\paragraph{DeepSeek-R1 (DeepSeek, 2025) \cite{deepseek_r1_2025}.}
DeepSeek-R1 dùng kiến trúc Mixture of Experts (MoE) với 671B tham số, trong đó 37B tham số được kích hoạt cho mỗi token. Điểm mới của mô hình là học tăng cường quy mô lớn cho suy luận mà không cần dữ liệu SFT ban đầu, nhờ đó mô hình tự kiểm tra và sửa lỗi trong chuỗi suy luận. DeepSeek-R1 đạt $78.17\%$ trên VM14K, cao nhất trong các mô hình được khảo sát \cite{vm14k_2025}, và $63.0\%$ trên OrthoBench \cite{orthopilot_2026}; bản chưng cất 70B đạt AUROC $0.837$ khi ước lượng độ bất định bằng Semantic Entropy \cite{mind_the_gap_2026}. Hạn chế chính là bản đầy đủ cần khoảng 960GB VRAM, vượt khả năng hạ tầng của phần lớn bệnh viện.

\paragraph{Qwen3 (Alibaba, 2025).}
Qwen3 là dải mô hình từ 0.6B đến 235B tham số, ngữ cảnh 128K--1M token, được huấn luyện trên dữ liệu đa ngôn ngữ và hỗ trợ cả chế độ trả lời trực tiếp lẫn chế độ suy luận (\emph{thinking}) trong cùng một mô hình. Qwen3-32B đạt $72.47\%$ trên VM14K, gần bằng GPT-4o \cite{vm14k_2025}; trên Spanish MIR, Qwen3-14B đạt $75.2\%$ và Qwen3-8B đạt $61.6\%$ \cite{spanish_mir_2026}. Tuy nhiên, các bản nhỏ còn yếu ở chuyên khoa: Qwen3-8B chỉ đạt $49\%$ ở nhóm câu hỏi chấn thương chỉnh hình, so với $80$--$85\%$ của các bản lớn hơn.

\paragraph{LLaMA-3.1 và LLaMA-4 (Meta, 2024--2025) \cite{llama3_2024}.}
LLaMA-3.1 dùng kiến trúc Dense với Grouped-Query Attention và ngữ cảnh 128K, còn LLaMA-4 Maverick chuyển sang kiến trúc MoE; cả hai được tiền huấn luyện đa ngôn ngữ, sau đó tinh chỉnh chỉ thị và căn chỉnh theo phản hồi con người. Trên VM14K, LLaMA-4 Maverick đạt $71.76\%$ và LLaMA-3.1-8B-Instruct đạt $48.73\%$; sau khi tinh chỉnh chỉ thị y khoa tiếng Việt, bản 8B tăng lên $59.2\%$, tức thêm 10.7 điểm \cite{vm14k_2025}. Bản 8B có tri thức y khoa tiếng Việt ban đầu thấp nên cần bổ sung bằng tinh chỉnh và truy xuất, nhưng bù lại có thể chạy trên một GPU 16--24GB.

Bảng~\ref{tab:general_llm_benchmark} tổng hợp kết quả của các mô hình đa dụng; GPT-4o được giữ làm mốc tham chiếu.

\begin{table}[H]
\begin{singlespace}
\centering
\fontsize{8pt}{10pt}\selectfont
\setlength{\tabcolsep}{3.5pt}
\renewcommand{\arraystretch}{1.1}
\begin{tabularx}{\textwidth}{|>{\raggedright\arraybackslash}X|>{\centering\arraybackslash}p{1.8cm}|>{\centering\arraybackslash}p{2.0cm}|>{\centering\arraybackslash}p{2.6cm}|>{\centering\arraybackslash}p{3.0cm}|}
\hline
\textbf{Mô hình} & \textbf{Loại} & \textbf{Quy mô} & \textbf{VM14K (Pass@1)} & \textbf{Spanish MIR (Accuracy)} \\
\hline
GPT-4o \cite{achiam2023gpt4} (tham chiếu) & Đóng & -- & $72.74\%$ & -- \\
\hline
DeepSeek-R1 \cite{deepseek_r1_2025} & Mở & MoE 671B & $\mathbf{78.17\%}$ & -- \\
\hline
Qwen3-32B & Mở & Dense 32B & $72.47\%$ & -- \\
\hline
Qwen3-14B & Mở & Dense 14B & -- & $75.2\%$ \\
\hline
Qwen3-8B & Mở & Dense 8B & -- & $61.6\%$ \\
\hline
LLaMA-4 Maverick & Mở & MoE & $71.76\%$ & -- \\
\hline
LLaMA-3.1-8B-Instruct \cite{llama3_2024} & Mở & Dense 8B & $48.73\%$ & -- \\
\hline
\end{tabularx}
\caption{Kết quả của các mô hình đa dụng trên VM14K và Spanish MIR}
\label{tab:general_llm_benchmark}
\end{singlespace}
\end{table}

%==========================================================
\subsection{Mô hình y khoa chuyên biệt}

Các mô hình chuyên biệt được chia thành năm nhóm theo cách đưa tri thức y khoa vào mô hình. Với nhóm 1 và nhóm 2, khóa luận chọn hai mô hình đại diện: mô hình đầu tiên của hướng tiếp cận và mô hình mới nhất đạt kết quả tốt nhất.

\subsubsection{Nhóm 1: Tiếp tục tiền huấn luyện trên y văn}
Ở hướng này, mô hình đa dụng được tiếp tục huấn luyện tự giám sát (\emph{continued pretraining}) trên hàng chục tỷ token y văn để học thuật ngữ và văn phong y khoa. Rủi ro chính là \emph{quên thảm họa} (\emph{catastrophic forgetting}): khi chỉ học trên y văn tiếng Anh, mô hình có thể mất năng lực ngôn ngữ khác như tiếng Việt.

\paragraph{PMC-LLaMA (Wu và cộng sự, 2023) \cite{wu2023pmc}.}
PMC-LLaMA được khởi tạo từ LLaMA 7B/13B và tiếp tục tiền huấn luyện trên 4.8 triệu bài báo toàn văn từ PubMed Central cùng sách giáo khoa y khoa, tổng cộng khoảng 75 tỷ token, sau đó được tinh chỉnh chỉ thị y khoa. Nhờ đó, độ chính xác trên MedQA-USMLE tăng từ $40.5\%$ lên $51.2\%$ so với LLaMA-7B gốc. Hạn chế là dữ liệu hoàn toàn bằng tiếng Anh, không có bệnh án và không có tri thức chuyên sâu về chấn thương chỉnh hình.

\paragraph{MEDITRON và Meditron3 (EPFL, 2023--2025) \cite{chen2023meditron}.}
Meditron-70B được xây trên LLaMA-2-70B và Meditron3-8B được xây trên LLaMA-3.1-8B, cả hai dùng tập GAP-Replay gồm 48.1 tỷ token từ PubMed, sách giáo khoa và hướng dẫn lâm sàng quốc tế. Để hạn chế quên kiến thức nền, nhóm tác giả trộn lại một phần dữ liệu đa dụng (\emph{replay}) trong quá trình tiền huấn luyện tiếp. Meditron-70B đạt $70.2\%$ trên MedQA, vượt GPT-3.5 và Med-PaLM. Tuy nhiên, trên VM14K, Meditron3-8B chỉ đạt $23.05\%$, thấp hơn một nửa so với mô hình gốc LLaMA-3.1-8B-Instruct ($48.73\%$) \cite{vm14k_2025}. Kết quả này cho thấy tiền huấn luyện tiếp chỉ bằng y văn tiếng Anh làm suy giảm mạnh năng lực tiếng Việt; ngoài ra bản 70B cần ít nhất hai GPU A100 80GB.

Tóm lại, tiền huấn luyện tiếp cải thiện điểm số tiếng Anh nhưng có thể làm hỏng năng lực tiếng Việt, nên không phù hợp làm mô hình nền cho bệnh án tiếng Việt.

\subsubsection{Nhóm 2: Tinh chỉnh chỉ thị trên dữ liệu y khoa}
Thay vì dùng văn bản thô, hướng này tinh chỉnh có giám sát (SFT) mô hình trên các bộ ba gồm chỉ thị, ca bệnh và phản hồi chuẩn, để mô hình học cách trả lời theo quy trình lâm sàng. Chi phí thấp hơn nhiều so với nhóm 1, nhưng có thể làm mô hình thận trọng quá mức và giảm độ chính xác trên bài toán suy luận phức tạp (\emph{alignment tax}).

\paragraph{HuatuoGPT-II (Chen và cộng sự, 2023) \cite{chen2023huatuogpt}.}
HuatuoGPT-II được xây trên Baichuan-2 7B/13B, dùng dữ liệu gồm sách giáo khoa y khoa, phác đồ và hội thoại bác sĩ--bệnh nhân thực tế. Thay vì tách riêng tiền huấn luyện và SFT, nhóm tác giả chuyển toàn bộ dữ liệu y khoa sang dạng chỉ thị--phản hồi và huấn luyện trong một giai đoạn duy nhất (\emph{one-stage}). Mô hình đạt kết quả tốt nhất tại thời điểm công bố trên các bài thi y khoa và đánh giá tư vấn tự do bằng tiếng Trung và tiếng Anh. Tuy vậy, mô hình thiên về tư vấn ngoại trú, chưa xử lý bệnh án có cấu trúc hay quy tắc chống chỉ định ngoại khoa.

\paragraph{MedGemma-27B (Google, 2025) \cite{spanish_mir_2026}.}
MedGemma-27B được tinh chỉnh từ Gemma-3-27B trên dữ liệu y khoa đa tác vụ gồm ca bệnh, câu hỏi thi chuyên khoa và hướng dẫn điều trị, tập trung vào phân tích ca bệnh và chẩn đoán phân biệt. Trên Spanish MIR, độ chính xác đạt $73.5\%$, \emph{thấp hơn} mô hình gốc ($75.8\%$); sai số hiệu chuẩn ECE giảm từ $0.231$ xuống $0.211$, còn AUROC gần như không đổi ($0.797$ so với $0.790$). Như vậy, tinh chỉnh y khoa chủ yếu làm mô hình bớt tự tin thái quá chứ không làm tăng độ chính xác chẩn đoán.

Tóm lại, tinh chỉnh chỉ thị giúp mô hình trả lời đúng định dạng lâm sàng nhưng không tự động tăng độ chính xác; cần so sánh mô hình trước và sau tinh chỉnh.

\subsubsection{Nhóm 3: Học tăng cường giảm ảo giác}
Hướng này dùng học tăng cường (\emph{reinforcement learning}) với mô hình phần thưởng y khoa để phạt các phát biểu sai sự thật và thưởng các lập luận có căn cứ.

\paragraph{Baichuan-M3-235B (Baichuan, 2026).}
Baichuan-M3 là mô hình MoE 235B tham số với ngữ cảnh 128K token. Mô hình kết hợp hai kỹ thuật: \emph{SPAR} chia quy trình khám thành các giai đoạn (hỏi bệnh, chỉ định cận lâm sàng, chẩn đoán, điều trị), mỗi giai đoạn có hàm thưởng riêng; còn học tăng cường có kiểm chứng sự thật (\emph{Fact-Aware}) phạt các phát biểu chẩn đoán không được phác đồ đã xác minh hỗ trợ. Trên HealthBench, tỷ lệ phát biểu bị bác bỏ giảm từ $4.68\%$ xuống $2.45\%$ và tỷ lệ phát biểu bất định giảm từ $3.64\%$ xuống $2.07\%$, trong khi điểm chất lượng chỉ giảm $1.1$ điểm. Dù vậy, tỷ lệ phát biểu sai vẫn khác 0, cho thấy học tăng cường không thay thế được nguồn tri thức bên ngoài; mô hình cũng cần cụm máy chủ lớn và chưa hỗ trợ tiếng Việt.

Tóm lại, giảm ảo giác bằng huấn luyện có hiệu quả nhưng chưa đủ; hệ thống vẫn cần kiểm chứng bằng chứng bên ngoài.

\subsubsection{Nhóm 4: Suy luận trên hồ sơ bệnh án điện tử}
Hồ sơ bệnh án điện tử (EHR) phân mảnh theo thời gian, chứa nhiều từ viết tắt và số liệu xét nghiệm. Hướng này huấn luyện mô hình suy luận trực tiếp trên tiến trình bệnh án thay vì trên câu hỏi thi.

\paragraph{EHR-R1 (SJTU, Shanghai AI Lab, Ant Group, 2025).}
EHR-R1 được xây trên Qwen với ba quy mô 1.7B, 8B và 72B, dùng tập chỉ thị EHR-Ins gồm 300.000 ca có suy luận và 3.5 triệu ca không suy luận trên 42 tác vụ lâm sàng, xây dựng từ MIMIC-IV. Mô hình được huấn luyện qua ba giai đoạn: thích ứng miền trên EHR thô, tinh chỉnh suy luận trên EHR-Ins, và căn chỉnh bằng học tăng cường theo kết cục lâm sàng thực tế. Bản 72B đạt F1 $0.6744$ trên EHR-Bench, cao hơn GPT-4o ($0.3155$) hơn 35 điểm, còn bản 8B vượt các mô hình đa dụng 70B trên các tác vụ dự đoán dọc thời gian. Hạn chế là mô hình chỉ được huấn luyện và đánh giá trên MIMIC-IV tiếng Anh và chưa tích hợp tri thức cơ xương khớp.

Tóm lại, suy luận trên bệnh án cần tổ chức thông tin theo cấu trúc và thời gian; quy mô mô hình hay ngữ cảnh dài không thay thế được yếu tố này.

\subsubsection{Nhóm 5: Mô hình và hệ thống chuyên khoa cơ xương khớp}
Chuyên khoa cơ xương khớp đòi hỏi kết hợp giải phẫu, cơ chế chấn thương, hệ phân loại gãy xương (AO/OTA, Garden, Schatzker) và quy tắc chống chỉ định. Nhóm này gồm hai tầng: mô hình lõi chuyên khoa và hệ thống agent dùng mô hình lõi kèm công cụ.

\paragraph{CHEESE-32B (Bệnh viện Ruijin, 2026) \cite{orthopilot_2026}.}
CHEESE-32B là mô hình Dense 32B, được huấn luyện chuyên khoa trên dữ liệu hồi cứu 21 năm của 183.472 bệnh nhân cơ xương khớp cùng các phác đồ phẫu thuật chỉnh hình. Trên OrthoBench (62 bệnh), mô hình đạt $68.9\%$, cao hơn GPT-5.1 ($67.2\%$) và DeepSeek-R1 ($63.0\%$); mô hình đạt $61.6\%$ trên 14 bệnh xương hiếm, $62.00\%$ ở chẩn đoán trước mổ và $71.00\%$ ở quản lý điều trị. Khi dùng độc lập, mô hình vẫn có nguy cơ ảo giác do tri thức cố định trong tham số, không đối chiếu được tài liệu gốc và chưa hỗ trợ tiếng Việt.

\paragraph{OrthoPilot (Bệnh viện Ruijin, 2026) \cite{orthopilot_2026}.}
OrthoPilot là hệ thống đa agent dùng CHEESE-32B làm mô hình lõi, kèm bộ công cụ tra cứu đồ thị tri thức chỉnh hình, truy xuất phác đồ và trích xuất bệnh án. Hệ thống vận hành theo vòng lặp ReAct: phân tích tổn thương, tra cứu phân loại gãy, duyệt đồ thị để tìm biến chứng, đối chiếu chống chỉ định, rồi lập kế hoạch điều trị. Trên OrthoBench, hệ thống đạt $75.4\%$ với 62 bệnh và $72.2\%$ với bệnh hiếm, $81.95\%$ ở chẩn đoán trước mổ và $85.13\%$ ở quản lý điều trị. Khi kiểm chứng ngoài tại 60 trung tâm (6.796 ca), OrthoPilot vượt DeepSeek-R1 $15.79$ điểm ở chẩn đoán; trong thử nghiệm ngẫu nhiên có đối chứng 8 tháng trên 8.240 bệnh nhân, số ca hoàn thành điều trị trên mỗi giường tăng từ $11.90$ lên $13.06$. Tuy nhiên, khi bỏ công cụ truy xuất và đồ thị tri thức, hiệu năng giảm về mức mô hình trần; hệ thống cũng chỉ được thử nghiệm tại Trung Quốc, chưa hỗ trợ tiếng Việt và chưa có cơ chế đo độ bất định.

Tóm lại, với cùng một mô hình lõi, việc bổ sung truy xuất bằng chứng và đồ thị tri thức giúp tăng $6.5$ đến $19.9$ điểm phần trăm. Đây là căn cứ trực tiếp cho kiến trúc RAG kết hợp đồ thị tri thức của khóa luận.

%==========================================================
\subsection{Bảng tổng hợp}

Bảng~\ref{tab:med_msk_comparison} tổng hợp các mô hình chuyên biệt theo cùng trình tự đã trình bày; Bảng~\ref{tab:llm_master_survey} so sánh các mô hình đại diện theo các tiêu chí triển khai trong bệnh viện.

\begin{table}[H]
\begin{singlespace}
\centering
\fontsize{7.5pt}{9.2pt}\selectfont
\setlength{\tabcolsep}{2pt}
\renewcommand{\arraystretch}{1.1}
\begin{tabularx}{\textwidth}{|>{\raggedright\arraybackslash}p{2.75cm}|>{\centering\arraybackslash}p{0.8cm}|>{\raggedright\arraybackslash}p{5.15cm}|>{\raggedright\arraybackslash}p{4.3cm}|>{\raggedright\arraybackslash}X|}
\hline
\textbf{Mô hình} & \textbf{Nhóm} & \textbf{Phương pháp} & \textbf{Kết quả tiêu biểu} & \textbf{Hạn chế chính} \\
\hline
PMC-LLaMA \cite{wu2023pmc} & 1 & Tiền huấn luyện tiếp trên y văn PubMed & MedQA: $40.5\% \rightarrow 51.2\%$ & Chỉ tiếng Anh, thiếu bệnh án \\
\hline
Meditron3-8B \cite{chen2023meditron} & 1 & Tiền huấn luyện tiếp kèm dữ liệu đa dụng & VM14K: $23.05\%$ (gốc $48.73\%$) & Suy giảm năng lực tiếng Việt \\
\hline
HuatuoGPT-II \cite{chen2023huatuogpt} & 2 & Tinh chỉnh chỉ thị trong một giai đoạn & Bài thi y khoa tiếng Trung, Anh & Chưa xử lý bệnh án cấu trúc \\
\hline
MedGemma-27B \cite{spanish_mir_2026} & 2 & Tinh chỉnh đa tác vụ từ Gemma-3-27B & Spanish MIR: $73.5\%$ (gốc $75.8\%$) & Không tăng độ chính xác \\
\hline
Baichuan-M3 & 3 & Học tăng cường có kiểm chứng sự thật & HealthBench, lỗi: $4.68\% \rightarrow 2.45\%$ & Quá lớn, chưa có tiếng Việt \\
\hline
EHR-R1-72B & 4 & Tinh chỉnh suy luận, học tăng cường & EHR-Bench: F1 $0.6744$ & Chỉ có MIMIC-IV tiếng Anh \\
\hline
CHEESE-32B \cite{orthopilot_2026} & 5 & Huấn luyện trên bệnh án cơ xương khớp & OrthoBench (62 bệnh): $68.9\%$ & Không đối chiếu tài liệu gốc \\
\hline
OrthoPilot \cite{orthopilot_2026} & 5 & Đa agent, tra cứu đồ thị và phác đồ & OrthoBench (62 bệnh): $75.4\%$ & Thiếu tiếng Việt, đo bất định \\
\hline
\end{tabularx}
\caption{Tổng hợp các mô hình y khoa chuyên biệt}
\label{tab:med_msk_comparison}
\end{singlespace}
\end{table}

\begin{table}[H]
\begin{singlespace}
\centering
\fontsize{7pt}{8.8pt}\selectfont
\setlength{\tabcolsep}{2.0pt}
\renewcommand{\arraystretch}{1.1}
\begin{tabularx}{\textwidth}{|>{\raggedright\arraybackslash}p{2.75cm}|>{\centering\arraybackslash}p{1.3cm}|>{\centering\arraybackslash}p{2.0cm}|>{\centering\arraybackslash}p{1.2cm}|>{\centering\arraybackslash}p{2.3cm}|>{\centering\arraybackslash}p{2.4cm}|>{\raggedright\arraybackslash}X|}
\hline
\textbf{Mô hình} & \textbf{Loại} & \textbf{Quy mô} & \textbf{VM14K} & \textbf{Chạy nội bộ} & \textbf{Đo bất định} & \textbf{Rủi ro chính} \\
\hline
GPT-4o \cite{achiam2023gpt4} & Đóng & -- & $72.74\%$ & \ding{55} (API) & \ding{55} & Dữ liệu bệnh nhân ra ngoài \\
\hline
DeepSeek-R1 \cite{deepseek_r1_2025} & Mở & MoE 671B & $78.17\%$ & \ding{55} ($\sim$960GB) & AUROC .837 (70B) & Quá nặng cho bệnh viện \\
\hline
Qwen3-32B & Mở & Dense 32B & $72.47\%$ & 2 GPU 24GB & AUROC .790 (14B) & Cần máy trạm 2 GPU \\
\hline
LLaMA-3.1-8B \cite{llama3_2024} & Mở & Dense 8B & $48.73\%$ & \ding{51} 1 GPU 16GB & Qua logit & Cần bổ sung tri thức tiếng Việt \\
\hline
Meditron3-8B \cite{chen2023meditron} & Y khoa & 8B & $23.05\%$ & \ding{51} 1 GPU 16GB & -- & Mất năng lực tiếng Việt \\
\hline
MedGemma-27B \cite{spanish_mir_2026} & Y khoa & Dense 27B & -- & 2 GPU 24GB & ECE .211 & Không tăng độ chính xác \\
\hline
Baichuan-M3-235B & Y khoa & MoE 235B & -- & \ding{55} (cụm GPU) & -- & Chưa có tiếng Việt \\
\hline
EHR-R1-72B & EHR & 72B & -- & $\ge 2$ A100 & -- & Chưa kiểm chứng tiếng Việt \\
\hline
CHEESE-32B \cite{orthopilot_2026} & MSK & Dense 32B & -- & 2 GPU 24GB & \ding{55} & Không đối chiếu nguồn \\
\hline
OrthoPilot \cite{orthopilot_2026} & Hệ thống & Agent & -- & Có thể & \ding{55} & Chưa có tiếng Việt \\
\hline
\textbf{Đề xuất} & \textbf{MSK} & \textbf{8B + QLoRA} & -- & \textbf{\ding{51} 1 GPU 16GB} & \textbf{\ding{51}} $u_{\mathrm{score}}$ & \textbf{Trích dẫn cấp câu, tiếng Việt} \\
\hline
\end{tabularx}
\caption{So sánh các mô hình đại diện theo tiêu chí triển khai lâm sàng}
\label{tab:llm_master_survey}
\end{singlespace}
\end{table}

%==========================================================
\subsection{Nhận xét chung}

Từ các kết quả trên, có thể rút ra năm nhận xét. Thứ nhất, mô hình đa dụng đa ngôn ngữ là điểm xuất phát tốt cho tiếng Việt: Qwen3-32B ($72.47\%$) và LLaMA-4 Maverick ($71.76\%$) gần bằng GPT-4o trên VM14K, trong khi Meditron3-8B chỉ đạt $23.05\%$; tinh chỉnh bằng dữ liệu y khoa tiếng Việt còn giúp bản 8B tăng hơn 10 điểm mà vẫn giữ năng lực ngôn ngữ.

Thứ hai, trả lời đúng không có nghĩa là giải thích đúng. Trên Spanish MIR, $13$--$38\%$ câu trả lời đúng đi kèm lý do không hợp lệ về mặt y khoa \cite{spanish_mir_2026}, nên cần đánh giá cả bằng chứng chứ không chỉ độ chính xác.

Thứ ba, các mô hình đều tự tin thái quá: độ tự tin trung bình từ $0.91$ đến $0.99$ trong khi độ chính xác chỉ từ $0.61$ đến $0.76$. Vì vậy, xác suất của mô hình chỉ nên dùng để xếp hạng và đặt ngưỡng từ chối, không dùng làm độ tin cậy chẩn đoán.

Thứ tư, huấn luyện chống ảo giác là chưa đủ. Baichuan-M3 giảm tỷ lệ phát biểu sai còn $2.45\%$ nhưng không về 0, do đó hệ thống cần nguồn tri thức ngoài và cơ chế neo bằng chứng.

Thứ năm, hệ thống có truy xuất và đồ thị tri thức vượt mô hình trần: so sánh OrthoPilot với CHEESE-32B cho thấy bổ sung công cụ truy xuất và đồ thị tăng $6.5$--$19.9$ điểm phần trăm.

%==========================================================
\subsection{Khoảng trống nghiên cứu ở Công đoạn 1}

\paragraph{Khoảng trống 1.1 -- Năng lực và bảo mật mâu thuẫn nhau.}
Các mô hình mạnh nhất như DeepSeek-R1 và GPT-4o hoặc phải gửi dữ liệu ra đám mây, hoặc cần hạ tầng khoảng 960GB VRAM. Ngược lại, mô hình nhỏ chạy được tại chỗ lại thiếu năng lực chuyên khoa nếu không được tinh chỉnh.

\paragraph{Khoảng trống 1.2 -- Thiếu tri thức cơ xương khớp tiếng Việt.}
Các mô hình y khoa mở được xây trên y văn tiếng Anh tổng quát, thiếu hệ phân loại gãy xương (AO/OTA, Garden, Schatzker), thuật ngữ giải phẫu tiếng Việt và các dấu hiệu cờ đỏ như gãy xương do u hay hoại tử chỏm xương đùi.

\paragraph{Khoảng trống 1.3 -- Thiếu cơ chế đo độ bất định.}
Hầu hết mô hình, kể cả CHEESE-32B, không tích hợp mô-đun đo độ bất định, nên vẫn đưa ra kết luận khi bệnh án thiếu xét nghiệm quan trọng.

\paragraph{Khoảng trống 1.4 -- Giải thích chưa kiểm chứng được.}
Chuỗi suy luận tự do có thể chỉ là hợp lý hóa sau khi đã có đáp án, và hiện chưa có cơ chế gán nguồn cho từng khẳng định lâm sàng.
